% Separate post-CP26 diagnostic note. No theorem promotion, hereditary PGL count, or review movement is asserted.
\subsection{The exact quantitative gap in transferring the projective atlas}

\paragraph{Status and question.}
This is an equation-by-equation audit of the constants in
Proposition~\ref{D:SNC-general-weighted-presentation} and
Theorem~\ref{D:SNC-general-projective-quadratic-count}.
Those results prove a quadratic count for each fixed complete
local surface. The question here is whether their constant on
a degree-$t$ punctured cover of a fixed base can be bounded by
$O_R(t)$. The answer does not follow from the currently stated
atlas estimates. This note identifies the unbounded factors;
it is not a counterexample to a hereditary projective theorem.

\paragraph{1. Geometric data that are now available uniformly.}
On the relative-minimal regular model of any such cover, the
all-complete geometric budgets and comparable-cycle theorem give
\[
 v+\sum_i(-E_i^2)+\sum_i p_a(E_i)+p_g=O_R(t),\qquad
 C\le Z\le\alpha C,                                      \tag{D.HP.1}
\]
where $C$ is the actual pullback of a fixed parameter cycle and
$Z=\alpha C-\beta R_g$ on that model. Moreover
\[
 -C^2,-C\cdot Z,-Z^2,|K\cdot Z|=O_R(t),\qquad
 h^0(\mathcal O_{bC})+h^1(\mathcal O_{bC})=O_{R,b}(t),       \tag{D.HP.2}
\]
and the corresponding statements hold on $bZ$.
Both $-C$ and $-Z$ are globally generated;
$H^1(Z,\mathcal O_Z(-jZ))=0$ for all $j\ge1$.
Canonical ordinary input bundles have the actual forward
Frobenius maps and inverse pole at most $pC$, and are generated
after twisting by $-Z$. These are the inputs of the proved
ordinary and ordinary-semilinear hereditary counts. They are
not yet statements about twisted bundles on a coefficient gerbe.

\paragraph{2. Pullback to an SNC refinement does not require a new Jacobian.}
Let $\rho:Y'\to Y$ be a sequence of point blowups and put
$C'=\rho^*C$, $Z'=\rho^*Z$. One must distinguish this definition
from recomputing $\alpha C'-\beta R_{g\rho}$. The latter need
not satisfy the old coefficient inequality after arbitrary
extra blowups. The former always satisfies
\[
 C'\le Z'\le\alpha C',\quad
 (C')^2=C^2,\quad C'\cdot Z'=C\cdot Z,\quad
 (Z')^2=Z^2,\quad K_{Y'}\cdot Z'=K_Y\cdot Z.              \tag{D.HP.3}
\]
For a point blowup, $\rho_*\mathcal O_{Y'}=\mathcal O_Y$
and $R^1\rho_*\mathcal O_{Y'}=0$; these follow from
\citet[Lemma 54.3.4(1),(2),(4), Tag 0AGS]{StacksHereditaryCostB}.
Composition, projection formula and the ideal sequence of
$bZ'$ therefore give
\[
 \rho_*\mathcal O_{bZ'}=\mathcal O_{bZ},\qquad
 R^1\rho_*\mathcal O_{bZ'}=0,\qquad
 H^1(Z',\mathcal O_{Z'}(-jZ'))
     =H^1(Z,\mathcal O_Z(-jZ))=0.                        \tag{D.HP.4}
\]
The same statement holds with $C$ in place of $Z$ and after
coefficient extension over $k$. Generation pulls back. Actual
ordinary input bundles and their arrows also pull back, retaining
the $pC'$ inverse bound and their original determinant list.
Their global sections are unchanged by projection formula, so
the target Hartogs property used by the generated-family Isom
argument is unchanged.

The pulled-back bundles need not be the canonical full pullbacks
of the same modules on $Y'$; that assertion is unnecessary and
is not made. Nor does (D.HP.4) prove that $\omega_{Z'}^{-1}$ is
generated. It preserves exactly the weak-cycle conditions used
by the new generated-source/fixed-determinant comparison.
New exceptional components can have
$C'\cdot E=Z'\cdot E=0$. A controlled number of SNC centres
is a separate question. Thus this observation removes a need
for a new Jacobian formula, not the componentwise positivity
or node-chart problem below.

\paragraph{3. Scalar classes and discrete inertia are not the rectangle obstruction.}
On an SNC resolution, (D.GG.2) immediately gives the simpler bound
\[
 |H^2(U,\mu_n)|\le n^{v+2\sum g_i+b_1},\qquad (n,p)=1,    \tag{D.HP.5}
\]
by using $|\operatorname{coker}(N\bmod n)|\le n^v$.
No determinant of the intersection matrix needs to be estimated
for this purpose. If SNC combinatorics are $O_R(t)$, the full
coefficient-gerbe classes therefore cost $q^{O_R(t)}$.

There is also a direct bound before SNC refinement. On a minimal
regular resolution with reduced exceptional curve $E$, remove
its finite singular/intersection set. Codimension-two purity
leaves $H^2(Y,\mu_n)$ unchanged, and codimension-one localization
on the remaining disjoint smooth components gives
\[
 |H^2(U,\mu_n)|\le
 n^{v+2\sum_i g(E_i^\nu)+b},                             \tag{D.HP.6}
\]
where $b$ is the total number of branches over the removed
points. Indeed proper base change identifies $H^2(Y,\mu_n)$
with $H^2(E,\mu_n)$, of order $n^v$, and the remaining term
is a subgroup of the $H^1$ of those punctured normalized curves.
The normalization defect at each removed point satisfies
$\delta_x\ge1$ and $\delta_x\ge b_x-1$, hence
$b\le2\delta(E)$. Since $E$ is connected,
\[
 \delta(E)=h^1(E,\mathcal O_E)-\sum_i g(E_i^\nu)+v-1,
 \qquad h^1(E,\mathcal O_E)\le p_g.
\]
Consequently (D.HP.6) is at most $n^{2p_g+3v}$, which is
$n^{O_R(t)}$ by (D.HP.1). The argument uses exactly the purity,
proper-base-change and normalized-curve inputs already checked
for (D.WD.5)--(D.WD.8); no assertion that the wild normal span
is pro-$p$ occurs. This is scalar coefficient control, not a
construction of a small twisted atlas.

For a fixed coefficient class, (D.GG.8) gives
\[
 g(R_i)+\#\operatorname{Marks}(\mathcal R_i)
       =O\bigl(m_i(1+g_i+\operatorname{valence}(i))\bigr).
                                                               \tag{D.HP.7}
\]
The last paragraph of the projective count bounds each node's
discrete scalar-weight-one type and its branch flags by
$\exp(O(fr\log(r+1)))$. Thus $O_R(t)$ nodes would cost
$\exp(O_R(ftr^2))$. The cyclic residue orders $m_i$ and their
least common multiple $M$ divide $r$. The rescaling $M/m_i$
already removes a field/period logarithm in (D.GP.11).
These observations do not remove the ordinary cycle coefficients.

\paragraph{4. The constants hidden by the fixed-resolution notation.}
Write $A=\sum_i a_iE_i$, $b_i=-E_i^2$,
$I_{ij}=E_i\cdot E_j$, and $c_i=-A\cdot E_i>0$ in the old
projective proof. Its componentwise argument requires strict
positivity of every $c_i$. Put
\[
 w_i=1+g_i+b_i+\operatorname{valence}(i),\qquad
 \epsilon_i(A)=a_i b_i+\sum_{j\ne i}a_j I_{ij}.          \tag{D.HP.8}
\]
The genus, marked points and minimal HN window in (D.GP.1),
(D.GP.5), (D.GG.8) have constants proportional to $w_i$.
The incomplete-cycle error in (D.TW.10), explicitly explained
just after that display, is bounded by a constant times
$m_i\epsilon_i(A)$. In (D.GP.6), one therefore chooses the
ordinary pole/vanishing depth using ratios of the form
\[
              \frac{w_i+\epsilon_i(A)}{c_i}.             \tag{D.HP.9}
\]
This is a sufficient-bound bookkeeping expression, not a
claimed necessary lower bound on every actual Hom space.
The existing theorem provides no cover-uniform bound for these
ratios. Substituting the new $Z$ is not a legal application of
that theorem when $-Z\cdot E_i=0$.

Even if a fixed precision $a$ were supplied by a different
argument, (D.GP.2) has $aa_i m_i$ component layers. If $\kappa_i$
denotes the actual per-layer Riemann--Roch/window constant
(the displayed proof allows it to depend on $w_i$ and the
incomplete-cycle errors), that coordinate count is
\[
 O\!\left(fr^2 a\sum_i a_i\kappa_i\right).              \tag{D.HP.10}
\]
This already need not be controlled by $v,\sum b_i,\sum g_i$
and the signed energy $-A^2$. In particular the proof has not
proved a bound for $\sum a_i\kappa_i$ from (D.HP.1)--(D.HP.2).

The node loss is completely explicit, with no hidden HN constant.
(D.GG.10), repeated as (D.TP.8) and used in (D.GP.8), gives
\[
 \#\text{ordinary coefficient slots at node }(i,j)
       \le a^2 a_i a_j r^2.                             \tag{D.HP.11}
\]
Summing the actual framed overlap/inverse data produces the
corresponding $O(fr^2a^2\sum_{\rm edges}a_i a_j)$ upper bound.
Root characters remove $m_i m_j$, but not $a_i a_j$.
This is a cost of that explicit presentation; it is not a
lower bound for every possible presentation of the moduli problem.

The weighted argument does not silently repair either loss.
In (D.GP.11), $J_i=aa_i m_i$ and the factorial ratio is
$((K_2m_i)^{J_i}/J_i!)^{b_i'}$, where
$b_i'=O(f\kappa_i r^2/m_i)$. Stirling therefore leaves terms
of the scale
\[
 fr^2a\sum_i a_i\kappa_i
       \bigl(1+\log^+(K_2/(aa_i))\bigr),                \tag{D.HP.12}
\]
in addition to the node/residual factors. The degree constant
$K_2$ itself depends on the largest retained ordinary orders and
on the fixed-complex constructions. The old conclusion uses
that $a_i$, $\kappa_i$, the number of nodes and $K_2$ are fixed.
No claim that they are all independent of a varying cover occurs
in the source theorem.

\paragraph{5. A precise numerical test of what energy cannot imply.}
Let the graph be the chain of length $2h-1$, with every square
$-2$, and let $z_i=\min(i,2h-i)$. Its positive matrix $N$ satisfies
\[
 Nz=2e_h,\quad z^tNz=2h,\quad
 \sum_i z_i=h^2,\quad
 \sum_{i=1}^{2h-2}z_i z_{i+1}=\frac23(h^3-h).           \tag{D.HP.13}
\]
Give only the middle vertex genus one. Its canonical vector is
$\kappa=2e_h$, and $\kappa^tN^{-1}\kappa=2h$.
Take $C=z$ and $Z=2z$. These cycles are anti-nef,
$C\le Z\le2C$, and their energies and canonical pairings are
$O(h)$. They even satisfy both all-subcycle numerical inequalities
used for a strong reduction cycle: if $B$ has middle coefficient
$u$, then
$2p_a(B)=2+2u-B^tNB\le2u$ because the $A$-chain quadratic
form is positive and even. Hence
$-Z\cdot B=4u$ and $-(K+Z)\cdot B=2u$ dominate $2p_a(B)$.
For $u=0$ the same formula has right side at most zero.
Nevertheless the node sum for either comparable cycle is cubic
in $h$, while its component coefficient sum is quadratic.

This is only a negative test of a proposed numerical implication.
It is not asserted to be a realized family of covers of one fixed
surface. It shows that the proved signed-energy and strong-cycle
inequalities, by themselves, cannot bound the unsigned rectangle
count in (D.HP.11). The actual cover geometry would need an
additional argument, or the atlas would need different coordinates.

\paragraph{6. What the new ordinary proof does, and what remains twisted.}
The ordinary proof uses an actual whole-cycle quotient
$[L^{-1}\to\mathcal O^{r+1}]$ and exact fixed-line derived Hom.
Its four blocks have total size $O_R(tr^2)$ on a fixed multiple
of the comparable cycle. It does not normalize the whole cycle
into independently framed components and rectangles. Raw finite
Frobenius effectivity and generated-source/fixed-determinant Isom
then remove the former coflat-stratum requirement.

The existing projective proof instead chooses its residual
Euler anchor of rank $M^2$ in (D.GP.3)--(D.GP.5), and an
algebraic generator $G_0$ in (D.GG.3)--(D.GG.4). Both are
explicitly used only for existence, minimization and algebraization;
neither rank is charged to the fixed-resolution atlas. Replacing
them by global quotient sources would require a new size proof.
In particular the existence of those anchors does not prove
that coarse Hom from them is generated on the same $Z$, that
its determinant degree on $bZ$ is $O(tr)$, or that its adjacent
Frobenius comparison has a cover-uniform Cartier pole bound.
Twisted determinants are not the ordinary characters enumerated
in Lemma~\ref{D:hereditary-ordinary-unit-determinants}.

For a bounded geometric period $M\le B$, rank $M^2$ is constant
in $t,r$. This removes one rank obstruction, but does not change
(D.HP.9)--(D.HP.11). A precise possible replacement input is a
choice of scalar-weight-$p^a$ anchors $A_a$, of rank bounded in
$B$, such that on the same comparable weak cycle:
\begin{enumerate}
\item for every relevant geometric twisted unit object, its
chosen coarse module $q_*\mathcal Hom(A_a,F_a)$ has an actual
whole-cycle generated presentation with rank $O_B(r)$ and
weighted determinant degree $O_{R,B}(tr)$;
\item the fixed coarse End algebras of the $A_a$ and the adjacent
Frobenius comparison bimodules have ranks bounded in $B$ and
actual fixed-complex cohomology/degree budgets $O_{R,B}(t)$;
\item their algebra actions, Morita conditions, arrows and inverse
relations admit exact all-ring whole-cycle equations with
$O_{R,B}(ftr^2)$ coordinates and degree bounded independently
of $t,r,f$; no ordinary matrix presentation with one independent
map for each of $O(t)$ algebra generators is substituted;
\item the retained generated data and determinant comparison
supply the actual specified-Artinian unit-Isom lifting required
for class openness on that twisted locus.
\end{enumerate}
These are prospective hypotheses, not proved assertions in this
note. The old fixed-period theorem's finitely many witness covers
are fixed only for one group; it gives no uniform degree bound
for splitters as the degree-$t$ subgroup varies. No such splitter
is used here. The next substantive task is the bounded-period
anchor/coarse-Hom input above, or a direct replacement for the
component/rectangle presentation. The existing arbitrary-SNC
projective theorem remains valid in its stated fixed-surface scope.
